\section{Introducci\'{o}n}

La microscopía electrónica de barrido (SEM) se ha consolidado como una herramienta esencial en el análisis de microestructuras en una amplia gama de disciplinas científicas y aplicaciones industriales. Esta técnica, que ha experimentado una evolución significativa desde su concepción, permite una observación detallada de las superficies a una escala microscópica, proporcionando perspectivas cruciales en campos tan diversos como la ciencia de materiales, la biomedicina y la ingeniería \parencite{Akhtar2018, SUN2023183, PARKER20223235}. La capacidad del SEM para revelar detalles intrincados de las superficies ha abierto nuevas vías para el análisis y la comprensión de propiedades materiales fundamentales. La exploración de las superficies a nivel microscópico ha sido un área de gran interés en la física y la ingeniería de materiales, dada su influencia crucial en las propiedades físicas y mecánicas de los materiales. En este contexto, la obtención y análisis de la rugosidad superficial emerge como un campo de estudio fundamental, no solo por su impacto en propiedades como la adhesión, la fricción y el desgaste, sino también por su papel en la innovación de materiales y tecnologías \parencite{SATO1982457, MONTAKHABI2022104737}. La presente investigación se centra en la obtención de parámetros de rugosidad a partir de imágenes estereoscópicas obtenidas mediante microscopía electrónica de barrido (SEM), una metodología que promete superar algunas de las limitaciones presentes en las técnicas convencionales mediante un enfoque no invasivo, versátil y eficiente.\\

